\subsection{Geometry of the morphospace}\label{subsec:geometry-of-the-morphospace}

The geometric analysis of the feature space reveals a distinctive interplay
between high-dimensional properties and evolutionary constraints.
The mean pairwise angle of $\frac{\pi}{2}$ aligns with the
property of high-dimensional spaces where random vectors tend to
be orthogonal and equidistant (\cref{fig:angles}).
Intra-taxa pairwise angles, although varied across taxa, are
smaller than the global mean value.
Since it is trained solely on species classification (flat labels),
a hierarchical taxonomic structure (Order/Family) automatically emerges in the high-dimensional feature space.
This demonstrates that feature vectors are semantically clustered,
a biological taxa occupying specific sub-manifolds and therefore exhibiting smaller
mean pairwise angle. This indicates that biological relationships are strongly encoded
in morphological features that models have to learn this structure when
searching for the optimal classification structure.

In the raw result of the clustering, many extinct species clustered together,
possibly due to their representation via skeletal images, artistic reconstructions, or
other non-biological patterns, being fundamentally different from extant species.
Furthermore, some of the recently described species or newly separated cryptic species grouped
together, being indistinguishable, which might reflect insufficient image data, leading to
underfitting of the model. All extinct species and some of the newly described or
newly split species, totally 249 species, were removed in subsequent analyses.

Although the clustering result is not perfectly align with current taxonomy,
these inconsistencies exhibits some historically noted morphological similarity.

In the Palaeognathae, four orders (Struthioniformes, Rheiformes, Casuariiformes, Apterygiformes)
of flightless birds form a distinct cluster, being driven by their unique body plans. In contrast,
Tinamiformes represents a notable exception, which is deeply within the Galliformes.
This topological incongruence highlights their convergent evolution. Tinamiformes and many Galliformes
have evolved highly convergently, exhibiting compact body shapes for ground-dwelling lifestyles,
and occupy similar niches \parencite{Glenny1946Systematic}. The model correctly identifies
this ecological and morphological similarity (\cref{fig:cluster}).

Furthermore, the inclusion of other visually quail-like terrestrial taxa like
Turnicidae (Buttonquails) \parencite{Rotthowe1998Evidence}, Menuridae \parencite{Winkler2020Lyrebirds},
Neomorphus \parencite{Haffer1977Neomorphus}
within the Galliformes cluster further validates that our model's feature space captures
morphological and ecological convergent adaptation and visual similarity. In the eyes of CNNs,
as long as you scratch for food and run around on the ground like a landfowl,
then you are a "landfowl" (\cref{fig:cluster}).

The clustering of Caprimulgiformes \textit{sensu lato}, Strigiformes, Accipitriformes, and Falconiformes
forms a visually similar group (\cref{fig:cluster}), with two layers of convergent evolution: Strigiformes and Caprimulgiformes \textit{s. l.}
share large eyes and cryptic plumage patterns for their nocturnal habits; and Strigiformes, Accipitriformes,
and Falconiformes share the raptorial facial characteristics like forward-facing eyes and hooked beaks
\parencite{Fidler2004Convergent, Wu2021Genomic}.

Other notable similarity exhibited in the clustering result includes
Hirundinidae and Apodi（Apodidae and Hemiprocnidae) \parencite{Norberg1986Insectivorous},
Sphenisciformes (Penguins), Alcidae and Procellariiformes \parencite{Cody1973Coexistence},
as well as Phaethontiformes, Suliformes, and Pelecanidae \parencite{Hedges1994Molecules}~\ldots~(\cref{fig:cluster})
It's demonstrated that the feature vectors' distribution in morphological space is not
arbitrary, but is constrained by objective phenotypic similarity. Considering that we did not
introduce \textit{a priori} taxonomic knowledge in the training process, it actually reinvented
the classical taxonomy from Linnaeus to Audubon. The model's taxonomy is based on phenotype,
yet with higher throughput of data than classical taxonomy and morphometrics.
This naturally creates a gap compared to molecular phylogenetic trees.
The gap represents the disconnection between visual disparity and genetic distance.

\section{Main drivers of diversification}\label{sec:main-drivers-of-diversification}
